% Ad M. Brutum 1.11 --- headnote
% ad-brutum-01-11, mid-to-late June 43 BC, from camp

\textit{M. Junius Brutus to Cicero, written from camp
in mid-to-late June 43 BC. The Perseus dateline reads
\textit{Scr.\ in castris ex.\ m.\ Iun., a.\ 711 (43)}
--- ``end of June 43 BC, in camp'' --- a month-precision
range; the meta entry imposes 15 June as a canonical
day-precision date, used in the parallel sidecar with
this discrepancy noted. Brutus is still in the Balkan
camp, after the campaign against C.\ Antonius and the
march along the Via Egnatia. The letter is a short
private commendation, but it characterises Brutus's
prose at its most warmly direct.}

\textit{The subject is ``veteran Antistius'' --- C.\
Antistius Vetus, the young senator who had supplied
Brutus with funds earlier in the year (see Ad Brutum
2.3.5) and who had refused, in Achaia, to give money to
Dolabella under the threat of his soldiers and cavalry.
Brutus tells the story of that refusal in a long
subordinate clause that builds suspense and lands on
its main verb: Antistius preferred to ``run any risk
you please from the ambushes of a brigand fully
equipped for every outrage'' rather than appear to have
been coerced into paying or willing to pay. He has now
joined Brutus in person and given twenty thousand
sesterces of his own money. Section 2 reports Brutus's
attempt to keep him in the camp as a commander rather
than let him depart for Rome to stand for the
praetorship, and his quiet success in dissuading him.
The closing request --- that Cicero ``love the veteran
and wish him as great a place as may be'' --- is the
note of personal warmth Brutus permits himself only
sparingly. The text in section 2 carries a small lacuna
(\textit{statuit id sibi * *}), here preserved with
asterisks where the manuscripts fail.}
